\documentclass{book}

\title{Quesiti parte orale Analisi Matematica T-A}
\author{Michele Mazza}
\date{November 2023}


% Pacchetti di Codifica e Lingua
\usepackage[utf8]{inputenc}
\usepackage[T1]{fontenc}
\usepackage[italian]{babel}

% Pacchetti Grafici e Diagrammi
\usepackage{graphicx}
\usepackage{tikz}
\usepackage{tikz-3dplot}
\usetikzlibrary{angles,quotes,arrows.meta, positioning}
\usepackage{pgfplots}
\pgfplotsset{compat=1.18}

% Pacchetti Matematici
\usepackage{amsmath, amssymb, amsthm}

% Pacchetti per Tabelle e Unità di Misura
\usepackage{geometry}
\usepackage{siunitx}
\usepackage{booktabs}
\usepackage{tabularx}
\usepackage{colortbl}

% Pacchetti per Layout e Formattazione
\usepackage{fancyhdr}
\usepackage{float}
\usepackage{microtype}

% Altri Pacchetti Utili
\usepackage{mdframed}
\usepackage{animate}

% Definizioni dei Colori
\definecolor{lightblue}{HTML}{9cd8fc}
\definecolor{myblue}{HTML}{139FE9}

% Definizione del comando \norm
\newcommand{\norm}[1]{\left\lVert#1\right\rVert}

%\font\nullfont=cmr10 

% Ambienti Personalizzati
\newmdtheoremenv[
    backgroundcolor=orange!10,
    linecolor=myblue,
    linewidth=0pt,
    leftmargin=0pt,
    rightmargin=0pt,
    innerleftmargin=10pt,
    innerrightmargin=10pt,
    innertopmargin=10pt,
    innerbottommargin=10pt
    nobreak
]{curiosity}{Curiosità}

\newmdtheoremenv[
    font=\normalfont\small,
    linewidth=0pt,
    innerleftmargin=10pt,
    innerrightmargin=10pt,
]{example}{Esempio}

\newmdtheoremenv[
  backgroundcolor=lightblue,
  linecolor=myblue,
  linewidth=1.5pt,
  innertopmargin=10pt,
  innerbottommargin=10pt,
  innerrightmargin=10pt,
  innerleftmargin=10pt
  nobreak
]{theorem}{Teorema}

\newmdtheoremenv[
  backgroundcolor=gray!10,
  linewidth=0pt,
  innertopmargin=10pt,
  innerbottommargin=10pt,
  innerrightmargin=10pt,
  innerleftmargin=10pt
  nobreak
]{definition}{Definizione}

\newmdtheoremenv[
  backgroundcolor=gray!10,
  linecolor=myblue,
  linewidth=1.5pt,
  innertopmargin=10pt,
  innerbottommargin=10pt,
  innerrightmargin=10pt,
  innerleftmargin=10pt
  nobreak
]{proposition}{Proposizione}

% Configurazioni di fancyhdr
\pagestyle{fancy}
\fancyhf{} % Pulisce tutti i campi
\fancyhead[LE,RO]{\thepage} % Numero di pagina
\fancyhead[RE]{\nouppercase{\leftmark}} % Titolo del capitolo
\fancyhead[LO]{\nouppercase{\rightmark}} % Titolo della sezione
\renewcommand{\headrulewidth}{0.5pt}
\renewcommand{\footrulewidth}{0pt}

\geometry{a4paper, left=20mm, right=20mm, top=20mm, bottom=20mm}

% Pacchetto Hyperref (inserito per ultimo)
\usepackage{hyperref}

\begin{document}

\frontmatter
\maketitle
\tableofcontents
\mainmatter
\chapter{La Retta Reale R}
Esploriamo le proprietà fondamentali della retta reale \( \mathbb{R} \), concentrandoci su relazioni d'ordine e di equivalenza, insiemi, estremi e limiti.

\subsection*{Relazioni d'Ordine e di Equivalenza}
Sia \( F \) un insieme non vuoto. Una relazione binaria \( A \) in \( F \) è definita come segue:
\begin{itemize}
  \item \textbf{Riflessiva:} \( A \) è riflessiva se \( aAa \) per ogni \( a \in F \).
  \item \textbf{Transitiva:} \( A \) è transitiva se da \( aAb \) e \( bAc \) segue \( aAc \) per ogni \( a, b, c \in F \).
  \item \textbf{Simmetrica:} \( A \) è simmetrica se da \( aAb \) segue \( bAa \) per ogni \( a, b \in F \).
  \item \textbf{Antisimmetrica:} \( A \) è antisimmetrica se da \( aAb \) e \( bAa \) segue \( a=b \) per ogni \( a, b \in F \).
\end{itemize}
Una \textbf{relazione di equivalenza} è riflessiva, simmetrica e transitiva. Una \textbf{relazione d'ordine parziale} è riflessiva, transitiva e antisimmetrica.

\subsection*{Insiemi Separati e Contigui}
\begin{itemize}
  \item Due insiemi si dicono \textbf{separati} se per ogni \( e \in F \) e \( g \in G \), \( e \neq g \).
  \item Sono \textbf{contigui} se sono separati e l'estremo superiore di uno coincide con l'estremo inferiore dell'altro.
\end{itemize}

\subsection*{Massimi e Minimi}
Per un sottoinsieme \( A \subseteq \mathbb{R} \):
\begin{itemize}
  \item \( A \) ha un \textbf{massimo} se esiste \( x_0 \in A \) tale che \( x_0 \geq x \) per ogni \( x \in A \).
  \item \( A \) ha un \textbf{minimo} se esiste \( x_0 \in A \) tale che \( x_0 \leq x \) per ogni \( x \in A \).
\end{itemize}

\subsection*{Estremo Superiore e Inferiore}
Un insieme \( A \) è:
\begin{itemize}
  \item \textbf{Limitato superiormente} se esiste \( K \in \mathbb{R} \) tale che \( K \geq x \) per ogni \( x \in A \). \( \sup(A) \) è il più piccolo \( K \) che soddisfa questa proprietà.
  \item \textbf{Limitato inferiormente} se esiste \( K \in \mathbb{R} \) tale che \( K \leq x \) per ogni \( x \in A \). \( \inf(A) \) è il più grande \( K \) che soddisfa questa proprietà.
\end{itemize}

\subsection*{Modulo di un Numero Reale}
Il \textbf{modulo} \( |x| \) di un numero reale \( x \) è definito come:
\[
  |x| =
  \begin{cases}
    x  & \text{se } x \geq 0, \\
    -x & \text{se } x < 0.
  \end{cases}
\]
Il modulo gode delle seguenti proprietà:
\begin{itemize}
  \item \( |x| = 0 \) se e solo se \( x = 0 \).
  \item \( |xy| = |x||y| \) per tutti \( x, y \in \mathbb{R} \).
  \item \( |x+y| \leq |x| + |y| \) (disuguaglianza triangolare).
\end{itemize}

\subsection*{Insiemi Limitati e Non Limitati}
Un insieme \( A \subset \mathbb{R} \) è \textbf{non limitato superiormente} se il suo estremo superiore è \( +\infty \) e \textbf{non limitato inferiormente} se il suo estremo inferiore è \( -\infty \).

\subsection*{Definizione di Limite}
Per una funzione \( f: A \to \mathbb{R} \) definita su un insieme \( A \subseteq \mathbb{R} \) non limitato superiormente, si ha:
\[
  \lim_{x \to +\infty} f(x) = L \in \mathbb{R}
\]
se e solo se per ogni \( \epsilon > 0 \) esiste \( M \) tale che per ogni \( x > M \) con \( x \in A \), \( |f(x) - L| < \epsilon \).

\chapter{Funzioni reali di variabile reale}
Sia \( f: A \to \mathbb{R} \) una funzione definita su un sottoinsieme \( A \) di \( \mathbb{R} \).

\subsection*{Limitatezza e Estremi di Funzioni}
Una funzione \( f \) si dice:
\begin{itemize}
  \item \textbf{Limitata superiormente} se esiste \( M \in \mathbb{R} \) tale che \( f(x) \leq M \) per ogni \( x \in A \). Il più piccolo di tali \( M \) è detto \(\sup f(A)\), l'estremo superiore di \( f \).
  \item \textbf{Limitata inferiormente} se esiste \( m \in \mathbb{R} \) tale che \( f(x) \geq m \) per ogni \( x \in A \). Il più grande di tali \( m \) è detto \(\inf f(A)\), l'estremo inferiore di \( f \).
  \item \textbf{Limitata} se è sia limitata superiormente che inferiormente.
\end{itemize}

\subsection*{Massimo e Minimo di Funzioni}
\begin{itemize}
  \item \( f \) ha un \textbf{massimo} in \( A \) se esiste \( x_0 \in A \) tale che \( f(x_0) \geq f(x) \) per ogni \( x \in A \). Tale \( f(x_0) \) è il massimo di \( f \), denotato come \(\max f(A)\).
  \item \( f \) ha un \textbf{minimo} in \( A \) se esiste \( x_0 \in A \) tale che \( f(x_0) \leq f(x) \) per ogni \( x \in A \). Tale \( f(x_0) \) è il minimo di \( f \), denotato come \(\min f(A)\).
\end{itemize}

\subsection*{Funzioni Monotone}
\begin{itemize}
  \item \( f \) si dice \textbf{monotona crescente} se per ogni \( x_1, x_2 \in A \) con \( x_1 < x_2 \), si ha \( f(x_1) \leq f(x_2) \).
  \item \( f \) si dice \textbf{monotona decrescente} se per ogni \( x_1, x_2 \in A \) con \( x_1 < x_2 \), si ha \( f(x_1) \geq f(x_2) \).
\end{itemize}

\subsection*{Limiti di Funzioni Monotone}
\begin{itemize}
  \item Se \( f \) è monotona crescente su \([a, b]\), allora \( \lim_{x \to b^-} f(x) = \sup f([a,b]) \).
  \item Se \( f \) è monotona decrescente su \([a, b]\), allora \( \lim_{x \to b^-} f(x) = \inf f([a,b]) \).
\end{itemize}

\subsection*{Teorema Generalizzato per Intervallo Infinito}
Per una funzione \( f: \mathbb{R} \to \mathbb{R} \) monotona:
\begin{itemize}
  \item Se \( f \) è monotona crescente nell'intervallo \( (-\infty, M) \), allora \( \lim_{x \to M^-} f(x) = \inf f((-\infty, M)) \).
  \item Se \( f \) è monotona decrescente nell'intervallo \( (M, +\infty) \), allora \( \lim_{x \to M^+} f(x) = \sup f((M, +\infty)) \).
\end{itemize}

\chapter{Il numero \( e \)}
Dimostriamo che la successione
\[ a_n = \left(1 + \frac{1}{n}\right)^n, \quad n = 1, 2, 3, \ldots \]
è crescente e limitata superiormente.

\subsection*{Crescita della Successione}
Consideriamo la successione \( a_n = \left(1 + \frac{1}{n}\right)^n \). Utilizziamo il Binomio di Newton che esprime \((1 + x)^n\) come una somma di termini binomiali:

\[ (1 + x)^n = \sum_{k=0}^{n} \binom{n}{k} x^k \]
dove
\[ \binom{n}{k} = \frac{n!}{k!(n-k)!}. \]

Applicando la disuguaglianza di Bernoulli, che afferma che per ogni \( a \geq 0 \) e \( n \in \mathbb{N} \):
\[ (1 + a)^n \geq 1 + na, \]
mostra che \( a_{n+1} \geq a_n \), ovvero la successione è crescente.

\subsection*{Limitatezza Superiore}
Osserviamo che \( a_n \) può essere espresso come
\[ a_n = \left(1 + \frac{1}{n}\right)^n = \sum_{k=0}^{n} \binom{n}{k} \left(\frac{1}{n}\right)^k = \sum_{k=0}^{n} \frac{1}{k!} \prod_{j=0}^{k-1} \left(\frac{n-j}{n}\right). \]
Notiamo che il prodotto \( \prod_{j=0}^{k-1} \left(\frac{n-j}{n}\right) \) è sempre minore o uguale a 1, quindi
\[ a_n \leq \sum_{k=0}^{n} \frac{1}{k!}, \]
e questa serie converge a \( e \), la base dei logaritmi naturali.

\subsection*{Conclusione}
La successione \( a_n \) è crescente e limitata superiormente. Il limite di questa successione, per \( n \to \infty \), è il numero \( e \), ovvero
\[ \lim_{n \to \infty} \left(1 + \frac{1}{n}\right)^n = e. \]
Quindi,
\[ \sup \left\{ \left(1 + \frac{1}{n}\right)^n : n \in \mathbb{N} \right\} = e. \]

\chapter*{Limiti di \( f: A \to \mathbb{R} \), con \( A \subseteq \mathbb{R} \)}

\subsection*{Punti di Accumulazione e Interni}
Un punto \( x_0 \in \mathbb{R} \) è detto di accumulazione di \( A \) se ogni intorno di \( x_0 \) contiene un punto di \( A \). Un punto è interno ad \( A \) se esiste un intorno completamente contenuto in \( A \).

\subsection*{Frontiera e Chiusura}
La frontiera di \( A \), denotata \( \partial A \), è composta da punti dove ogni intorno interseca sia \( A \) che il suo complemento. La chiusura di \( A \), \( \overline{A} \), include tutti i punti aderenti ad \( A \).

\subsection*{Limite di una Restrizione}
Per \( f: A \to \mathbb{R} \) e \( x_0 \) punto di accumulazione, se \( \lim_{x \to x_0} f(x) = L \), allora per ogni sottoinsieme \( B \subseteq A \), \( \lim_{x \to x_0, x \in B} f(x) = L \).

\subsection*{Limiti Destro e Sinistro}
Il limite destro di \( f \) in \( x_0 \) è definito su \( A^+ = A \cap [x_0, +\infty) \) e analogamente per il limite sinistro su \( A^- = A \cap (-\infty, x_0] \).

\subsection*{Teoremi sui Limiti}
I limiti sono locali e implicano la limitatezza di \( f \) in un intorno di \( x_0 \). Conservano il segno e seguono le proprietà delle funzioni composte.

\subsection*{Confronto di Infinitesimi e Infiniti}
Il confronto tra infinitesimi o infiniti si basa sul rapporto tra le funzioni. Due funzioni sono infinitesimi dello stesso ordine se il loro rapporto tende a un limite finito e non nullo.

\section*{Dimostrazione del limite \( \lim_{x \to 0}\frac{\sin x}{x} = 1 \)}
Per dimostrare \( \lim_{x \to 0}\frac{\sin x}{x} = 1 \), consideriamo la disuguaglianza \( \cos x < \frac{\sin x}{x} < 1 \) per \( 0 < x < \frac{\pi}{2} \). La continuità del coseno in \( 0 \) e il Teorema del Confronto confermano il limite.

\end{document}